\section{引言}

MoE 以稀疏路由减少单 token 激活的专家数，但低内存机器仍必须在映射权重、页缓存、KV cache 和临时缓冲之间做出及时而安全的取舍。已有端侧工作说明，权重卸载、内存保留和异步数据移动是重要的设计空间：FlexInfer 探索了张量级卸载与保留\cite{flexinfer2025}，PowerInfer-2 面向智能手机讨论了端侧大模型执行\cite{powerinfer2_2024}，MoE-Prism 则表明模型与系统协同能够改变专家服务的资源弹性\cite{moeprism2025}。这些工作构成问题和设计动机；它们并不证明 SLIM-ARC 的实现或性能。

决赛阶段的问题更具体：当 Router 的短期预测与实际选择不一致，如何撤回不再有用的页建议；当内存压力升高时，如何只为有足够证据的专家保留有限预算；以及如何证明建议目标在并发和模型析构下仍然有效。把这些问题仅当作启发式调参，会把错误的地址、重复工作或悬空后台任务带入推理路径。

SLIM-ARC 因而将优化表述为一个受约束的闭环：\emph{预测 $\rightarrow$ 准入 $\rightarrow$ 预取 $\rightarrow$ 观测 $\rightarrow$ 回收}。本文贡献是：
\begin{itemize}
  \item 给出一套仅对完整内部页操作的错误专家回收计划，并将非法布局、非法 ID、跳过字节和建议失败显式计数；
  \item 给出压力、错误率和热度共同约束的专家驻留策略，所有新行为保持 opt-in，未通过性能门槛时不成为默认路径；
  \item 用模型所有权、租约和有界工作队列约束生命周期与并发，并通过严格运行时指标将机制行为连接到每次实验；
  \item 规定单一的 Mac 受限环境实验协议和结果导入边界。正式结果只能由机器生成的 JSON 及其原始运行证据支持。
\end{itemize}
